\section*{Sets and Indices}

\[
P = \{\text{A1},\text{A2},\text{A3}\}
\]

Index:
\[
p \in P
\]

\section*{Parameters}

For each product \(p \in P\):
\begin{itemize}
    \item \(D_p\): maximum demand (from \texttt{Maximum\_Demand})
    \item \(s_p\): selling price per unit (from \texttt{Selling\_Price})
    \item \(c_p\): production cost per unit (from \texttt{Production\_Cost})
    \item \(q_p\): production quota in units per operating day (from \texttt{Production\_Quota})
    \item \(F_p\): activation cost (from \texttt{Activation\_Cost})
    \item \(B_p\): minimum batch size (from \texttt{Minimum\_Batch})
\end{itemize}

Global parameters:
\begin{itemize}
    \item \(T\): available production days in the month (code/data: \texttt{production\_days})
    \item \(\theta^{W}\): A1 watch threshold (code/data: \texttt{a1\_watch\_threshold})
    \item \(\theta^{D}\): A1 deep-service threshold (code/data: \texttt{a1\_deep\_service\_threshold})
    \item \(\delta^{W}\): operating-day loss if the watch threshold is crossed (code/data: \texttt{maintenance\_penalty\_watch})
    \item \(\delta^{D}\): additional operating-day loss if the deep-service threshold is crossed (code/data: \texttt{maintenance\_penalty\_deep})
    \item \(K^{W}\): quality-release fee charged if the watch threshold is crossed (code/data: \texttt{a1\_quality\_release\_fee})
    \item \(K^{D}\): deep-service fee charged if the deep-service threshold is crossed (code/data: \texttt{a1\_deep\_service\_fee})
    \item \(\pi\): priority price uplift for A1 units above the watch threshold (code/data: \texttt{a1\_priority\_price\_lift})
\end{itemize}

Auxiliary constant used in the code:
\[
M = D_{\text{A1}}
\]

\section*{Decision Variables}

\begin{itemize}
    \item \(x_p \ge 0\) (code: \texttt{x[p]}): production quantity of product \(p\)
    \item \(y_p \in \{0,1\}\) (code: \texttt{y[p]} / \texttt{open[p]}): product activation indicator
    \item \(w \in \{0,1\}\) (code: \texttt{watch} / \texttt{a1\_watch}): indicator for A1 crossing the watch threshold
    \item \(d \in \{0,1\}\) (code: \texttt{deep} / \texttt{a1\_deep}): indicator for A1 crossing the deep-service threshold
    \item \(e \ge 0\) (code: \texttt{excess} / \texttt{a1\_priority\_excess}): A1 units eligible for the priority uplift
\end{itemize}

\section*{Objective}

Maximize total operating contribution plus A1 priority uplift, minus activation costs and threshold-triggered charges:
\[
\max \;
\sum_{p \in P} (s_p-c_p)x_p
\;+\;
\pi e
\;-\;
\sum_{p \in P} F_p y_p
\;-\;
K^{W} w
\;-\;
K^{D} d
\]

\section*{Constraints}

All products are forced open in the implemented model:
\[
y_p = 1 \qquad \forall p \in P
\]

Demand and minimum-batch bounds:
\[
x_p \le D_p\, y_p \qquad \forall p \in P
\]
\[
x_p \ge B_p\, y_p \qquad \forall p \in P
\]

A1 watch-threshold logic:
\[
x_{\text{A1}} - \theta^{W} \le (M-\theta^{W})\, w
\]

A1 deep-service logic:
\[
x_{\text{A1}} - \theta^{D} \le (M-\theta^{D})\, d
\]

Deep service implies watch:
\[
d \le w
\]

Priority-uplift excess for A1:
\[
e \ge x_{\text{A1}} - \theta^{W}
\]
\[
e \le x_{\text{A1}} - \theta^{W} + M(1-w)
\]
\[
e \le (M-\theta^{W})\, w
\]

Available operating days, including threshold-triggered day losses:
\[
\sum_{p \in P} \frac{x_p}{q_p}
\;\le\;
T - \delta^{W} w - \delta^{D} d
\]

\section*{Notes on Implementation}

The formulation above matches the model implemented in \texttt{current\_heuristic.py}.

In particular:
\begin{itemize}
    \item The binary activation variables \(y_p\) are present, but the code imposes \(y_p=1\) for every product, so all three product lines must be active.
    \item The watch and deep-service rules are modeled as gates, not proportional surcharges: if A1 output exceeds \(\theta^W\), then \(w\) must become 1; if it exceeds \(\theta^D\), then \(d\) must become 1. These binaries trigger fixed fees and reduce available production days.
    \item The priority uplift applies only to A1 units above the watch threshold. Variable \(e\) is linearized so that, in combination with profit maximization, it equals
    \[
    e = \max\{0,\, x_{\text{A1}}-\theta^W\}
    \]
    with the additional restriction that positive \(e\) is only possible when \(w=1\).
    \item The monthly capacity constraint uses production days consumed as \(\sum_p x_p/q_p\), then subtracts \(\delta^W\) and \(\delta^D\) from available days when the corresponding A1 thresholds are crossed.
    \item The code solves this as a mixed-integer linear program using \texttt{GUROBI\_CMD}.
\end{itemize}
